\documentclass[a4paper]{article}
\usepackage{ctex}
\usepackage[affil-it]{authblk}
\usepackage[backend=bibtex,style=numeric]{biblatex}
\usepackage{amsmath,amsthm,amssymb,amsfonts}
\usepackage{bm}

\usepackage{geometry}
\geometry{margin=1.5cm, vmargin={0pt,1cm}}
\setlength{\topmargin}{-1cm}
\setlength{\paperheight}{29.7cm}
\setlength{\textheight}{25.3cm}

\begin{document}
% =================================================
\title{计算电磁学 2027 秋冬作业一}

\author{曾申昊 12631040
  \thanks{邮箱: \texttt{3097714673@qq.com}
                            \texttt{或} \texttt{12631040@zju.edu.cn}}}
\affil{浙江大学}

\date{更新时间: \today}

\maketitle

% ============================================
\section*{Problem 2.2}
 
Determine the distribution of electrostatic potential inside 
the conducting rectangular boxes with cross sections shown in 
Fig. 2.18.

\begin{enumerate}
    \item[(a)] box 1
        \begin{gather*}
            \nabla^2 U = 0, \quad -\frac{a}{2} < x < \frac{a}{2}, \quad -\frac{b}{2} < y < \frac{b}{2} \\
            U\left(-\frac{a}{2}, y\right) = V_0 = U\left(\frac{a}{2}, y\right),\\
            U\left(x, -\frac{b}{2}\right) = 0 = U\left(x, \frac{b}{2}\right)
        \end{gather*}
        使用分离变量法求解方程，将 \(U(x, y) = X(x)Y(y)\) 代入 Laplace equation。
        \begin{gather*}
            \nabla^2 (XY) = 0 \\
            X''Y + XY'' = 0\\
            \frac{X''}{X} + \frac{Y''}{Y} = 0
        \end{gather*}
        等价于
        \begin{equation*}
            \begin{cases}
                X'' - \beta^2 X = 0 \\
                Y'' + \beta^2 Y = 0
            \end{cases}
        \end{equation*}
        令 $\tilde{y}=y+b/2\in(0,b)$，并利用上下边界的齐次条件可得 $\beta_n = \frac{n\pi}{b}$ 以及：
        \begin{equation*}
            U(x,y)=\sum_{n=1}^{\infty}X_n(x)
            \sin\!\left(\frac{n\pi\tilde{y}}{b}\right)
        \end{equation*}
        代入 Laplace 方程可得
        \begin{equation*}
            X_n''-\left(\frac{n\pi}{b}\right)^2X_n=0.
        \end{equation*}
        两个竖直边界上的电势相同，因此取偶解并由边界条件确定
        \begin{equation*}
            X_n(x)=A_n\frac{\cosh\left(\frac{n\pi x}{b}\right)}{\cosh\left(\frac{n\pi a}{2b}\right)}.
        \end{equation*}
        在 $x=\pm \frac{a}{2}$ 处，$U=V_0$，故需将常数展开为正弦级数：
        \begin{equation*}
            V_0=\sum_{n=1}^{\infty}A_n
            \sin\!\left(\frac{n\pi\tilde{y}}{b}\right),\qquad
            A_n=\frac{2V_0}{b}\int_0^b\sin\!\left(\frac{n\pi\tilde{y}}{b}\right)\,d\tilde{y}
            =\begin{cases}\displaystyle \frac{4V_0}{n\pi},&n\text{ odd},\\[4pt]0,&n\text{ even}.
            \end{cases}
        \end{equation*}
        因而 box 1 内部的电势分布为
        \begin{equation*}
            U(x,y)=\frac{4V_0}{\pi}\sum_{m=0}^{\infty}\frac{1}{2m+1}
            \frac{\cosh\!\left(\dfrac{(2m+1)\pi x}{b}\right)}
                {\cosh\!\left(\dfrac{(2m+1)\pi a}{2b}\right)}
            \sin\!\left[\frac{(2m+1)\pi}{b}\left(y+\frac b2\right)\right]
        \end{equation*}
    \item[(b)] box 2
        \begin{gather*}
            \nabla^2 U = 0, \quad x>0, \quad 0 < y < b \\
            U\left(0, y\right) = V_0,\\
            U\left(x, 0\right) = 0 = U\left(x, b\right)
        \end{gather*}
        同上一小问可得 $\beta_n = \frac{n\pi}{b}$ 以及：
        \begin{equation*}
            U(x,y)=\sum_{n=1}^{\infty}X_n(x)
            \sin\!\left(\frac{n\pi y}{b}\right)
        \end{equation*}
        为了使得正无穷远处的值有限，因此取指数衰减解：
        \begin{equation*}
            X_n(x)=A_n e^{-\frac{n\pi x}{b}}
        \end{equation*}
        在 $x=0$ 处，$U=V_0$，故需将常数展开为正弦级数：
        \begin{equation*}
            V_0=\sum_{n=1}^{\infty}A_n
            \sin\!\left(\frac{n\pi y}{b}\right),\qquad
            A_n=\frac{2V_0}{b}\int_0^b\sin\!\left(\frac{n\pi y}{b}\right)\,dy
            =\begin{cases}\displaystyle \frac{4V_0}{n\pi},&n\text{ odd},\\[4pt]0,&n\text{ even}.
            \end{cases}
        \end{equation*}
        因而 box 1 内部的电势分布为
        \begin{equation*}
            U(x,y)=\frac{4V_0}{\pi}\sum_{m=0}^{\infty}\frac{1}{2m+1}
            e^{-\frac{(2m+1)\pi x}{b}}
            \sin\!\left[\frac{(2m+1)\pi}{b}y\right]
        \end{equation*}
    \end{enumerate}

    \section*{Problem 2.4}

    Find the solution \(U\) of:

    \begin{enumerate}
    \item[(a)] Laplace equation
        \begin{align*}
            \nabla^2 U &= 0, \quad 0 < x, y < \pi \\
            U_x(0, y) &= 0 = U_x(\pi, y), \quad U(x, 0) = 0, \\
            U(x, \pi) &= x, \quad 0 < x < \pi
        \end{align*}
        使用分离变量法求解方程，将 \(U(x, y) = X(x)Y(y)\) 代入 Laplace equation。
        \begin{gather*}
            \nabla^2 (XY) = 0 \\
            X''Y + XY'' = 0\\
            \frac{X''}{X} + \frac{Y''}{Y} = 0
        \end{gather*}
        等价于
        \begin{equation*}
            \begin{cases}
                X'' + \beta^2 X = 0 \\
                Y'' - \beta^2 Y = 0
            \end{cases}
        \end{equation*}
        利用上下边界的导数齐次条件可得 $\beta_n = n$ 以及：
        \begin{equation*}
            U(x,y)=\sum_{n=1}^{\infty}\cos\!\left(n x\right)
            Y_n(y)
        \end{equation*}
        代入 Laplace 方程可得
        \begin{equation*}
            Y_n''-\left(\frac{n\pi}{b}\right)^2Y_n=0.
        \end{equation*}
        由边界条件 $U(x,0)=0$ 确定 
        \begin{equation*}
            Y_n(y) = B_n \sinh\left( n y\right)
        \end{equation*}
        在 $y=\pi$ 处，$U=x$，故需将 $x$ 展开：
        \begin{equation*}
            x = \sum_{n=1}^{\infty} B_n \cos(nx) \sinh\left(n \pi\right)
        \end{equation*}
        \begin{equation*}
            \begin{aligned}
                B_n &= \frac{2}{\pi\sinh(n \pi)}\int_0^\pi x \cos(nx) \, dx \\
                &= \frac{2}{\pi\sinh(n \pi)}\left[\frac{\pi}{n}\sin(n \pi) + \frac{1}{n^2}\cos(n \pi) - \frac{1}{n^2}\right] \\
                &= \frac{2}{\pi\sinh(n \pi)}\left[\frac{1}{n^2}(-1)^n - \frac{1}{n^2}\right] \\
                &= \frac{2((-1)^n - 1)}{n^2\pi\sinh(n \pi)}
            \end{aligned}
        \end{equation*}
        因而 $U(x,y)$ 为
        \begin{equation*}
            \begin{aligned}
                U(x,y) &= \sum_{n=1}^{\infty} \frac{2((-1)^n - 1)}{n^2\pi\sinh(n \pi)} \cos(nx) \sinh(n y)\\
                &= \sum_{m=0}^{\infty} \frac{-4}{(2m+1)^2\pi\sinh((2m+1) \pi)} \cos((2m+1)x) \sinh((2m+1) y)
            \end{aligned}
        \end{equation*}
    \item[(b)] Heat equation
        \begin{align*}
            kU_{xx} &= U_t, \quad 0 \leq x \leq 1, \; t > 0 \\
            U(0, t) &= 0, \; t > 0, \quad U(1, t) = 1, \; t > 0 \\
            U(x, 0) &= 0, \quad 0 \leq x \leq 1
        \end{align*}
        令 $U(x, t) = x + \tilde{U}(x,t)$ 可得
        \begin{align*}
            k\tilde{U}_{xx} &= \tilde{U}_t, \quad 0 \leq x \leq 1, \; t > 0 \\
            \tilde{U}(0, t) &= 0, \; t > 0, \quad \tilde{U}(1, t) = 0, \; t > 0 \\
            \tilde{U}(x, 0) &= -x, \quad 0 \leq x \leq 1
        \end{align*}
        使用分离变量法求解方程，将 \(\tilde{U}(x, t) = X(x)T(t)\) 代入 heat equation。
        \begin{gather*}
            kX''T = XT'\\
            k\frac{X''}{X} = \frac{T'}{T}
        \end{gather*}
        等价于
        \begin{equation*}
            \begin{cases}
                X'' + \frac{\lambda}{k} X = 0 \\
                T' + \lambda T = 0
            \end{cases}
        \end{equation*}
        令 $\lambda = kn^2\pi^2$ 得到
        \begin{equation*}
            \begin{cases}
                X'' + n^2\pi^2 X = 0 \\
                T' + kn^2\pi^2 T = 0
            \end{cases}
        \end{equation*}
        由边界条件易得 $X_n(x) = \sin(n\pi x)$，
        解的整体形式可以写为
        \begin{equation*}
            \tilde{U}(x, t) = \sum_{n=1}^{\infty} A_n \sin(n\pi x) e^{-kn^2\pi^2 t}
        \end{equation*}
        由初始条件 $\tilde{U}(x, 0) = -x$ 可得
        \begin{equation*}
            -x = \sum_{n=1}^{\infty} A_n \sin(n\pi x)
        \end{equation*}
        因而
        \begin{equation*}
            A_n = 2\int_0^1 (-x) \sin(n\pi x) \, dx
            = \frac{2}{n\pi}(-1)^n
        \end{equation*}
        最终得到
        \begin{equation*}
            U(x, t) = x + \sum_{n=1}^{\infty} \frac{2}{n\pi}(-1)^n \sin(n\pi x) e^{-kn^2\pi^2 t}
        \end{equation*}
    \item[(c)] Wave equation
        \begin{align*}
            a^2 U_{xx} &= U_{tt}, \quad 0 \leq x \leq 1, \; t > 0 \\
            U(0, t) &= 0 = U(1, t), \; t > 0 \\
            U(x, 0) &= 0, \quad U_t(x, 0) = x
        \end{align*}
        使用分离变量法求解方程，将 \(U(x, y) = X(x)T(t)\) 代入 wave equation。
        \begin{gather*}
            a^2 X''T = XT''\\
            a^2\frac{X''}{X} = \frac{T''}{T}
        \end{gather*}
        等价于
        \begin{equation*}
            \begin{cases}
                X'' + \frac{\beta^2}{a^2} X = 0 \\
                T'' + \beta^2 T = 0
            \end{cases}
        \end{equation*}
        并利用上下边界的齐次条件可得 $\beta_n = na\pi$ 以及：
        \begin{equation*}
            U(x,t)=\sum_{n=1}^{\infty}\sin\left(n\pi x\right)T_n(t)
        \end{equation*}
        代入可得
        \begin{equation*}
            T'' + (na\pi)^2 T = 0
        \end{equation*}
        由边界条件 $U(x,0) = 0$ 确定
        \begin{equation*}
            T(t) = A_n\sin\left(na\pi t\right)
        \end{equation*}
        由边界条件 $U_t(x,0) = x$ 可得
        \begin{equation*}
            x = \sum_{n=1}^{\infty} \sin\left(n\pi x\right) A_n (na\pi)
        \end{equation*}
        \begin{equation*}
            A_n = \frac{2}{na\pi}\int_0^1 x \sin\left(n\pi x\right) \, dx
            = \frac{-2}{an^2\pi^2}(-1)^n
        \end{equation*}
        因而
        \begin{equation*}
            U(x,t) = \sum_{n=1}^{\infty} \frac{-2}{an^2\pi^2}(-1)^n \sin\left(n\pi x\right) \sin\left(na\pi t\right)
        \end{equation*}
    \end{enumerate}

    \section*{Problem 2.5}

    Find the solution \(\Phi\) of:

    \begin{enumerate}
    \item[(a)] Laplace equation
        \begin{align*}
            \nabla^2 \Phi &= 0, \qquad \rho \geq 1, \, 0 < \phi < \pi \\
            \Phi(1, \phi) &= \sin \phi, \qquad \Phi(\rho, 0) = \Phi(\rho, \pi) = 0
        \end{align*}
        在极坐标系中，Laplace 方程为
        \begin{equation*}
            \frac{\partial^2\Phi}{\partial \rho^2}
            +\frac{1}{\rho}\frac{\partial\Phi}{\partial\rho}
            +\frac{1}{\rho^2}\frac{\partial^2\Phi}{\partial\phi^2}=0.
        \end{equation*}
        令 $\Phi(\rho,\phi)=R(\rho) F(\phi)$，代入并乘以
        $\rho^2/(R F)$，得
        \begin{equation*}
            \frac{\rho^2R''+\rho R'}{R}
            =-\frac{ F''}{ F}=\lambda.
        \end{equation*}
        因此分离后的方程为
        \begin{equation*}
            \begin{cases}
                F''+\lambda F=0,\\
                \rho^2R''+\rho R'-\lambda R=0.
            \end{cases}
        \end{equation*}
        由 $ F(0)= F(\pi)=0$ 可得
        \begin{equation*}
            \lambda_n=n^2,\qquad  F_n(\phi)=\sin(n\phi),\qquad
            n=1,2,\ldots
        \end{equation*}
        相应的径向方程为 Euler 方程，其解为
        \begin{equation*}
            R_n(\rho)=A_n\rho^n+B_n\rho^{-n}.
        \end{equation*}
        对外部区域要求电势在 $\rho\to\infty$ 时有界，
        故只保留衰减项 $\rho^{-n}$，于是
        \begin{equation*}
            \Phi(\rho,\phi)=\sum_{n=1}^{\infty}B_n\rho^{-n}\sin(n\phi).
        \end{equation*}
        在 $\rho=1$ 处利用边界条件
        \begin{equation*}
            \sin\phi=\sum_{n=1}^{\infty}B_n\sin(n\phi),
        \end{equation*}
        由正弦函数的正交性可得 $B_1=1$，其余 $B_n=0$。所以
        \begin{equation*}
            \Phi(\rho,\phi)=\frac{\sin\phi}{\rho}
        \end{equation*}
    \item[(b)] Laplace equation
        \begin{align*}
            \nabla^2 \Phi &= 0, \\
            \Phi(\rho, \phi, 0) &= 0 = \Phi(\rho, \phi, L), \Phi(a, \phi, z) = 1
        \end{align*}
        展开 Laplace 方程可得
        \begin{align*}
            \frac{1}{\rho} \frac{\partial}{\partial \rho} \left(\rho \frac{\partial \Phi}{\partial \rho}\right) + \frac{1}{\rho^2} \frac{\partial^2 \Phi}{\partial \phi^2} + \frac{\partial^2 \Phi}{\partial z^2}&= 0
        \end{align*}
        令$\Phi(\rho, \phi, z) = R(\rho) F(\phi) Z(z)$并代入上述方程，得
        \begin{align*}
            \frac{\rho}{R} \frac{d}{d\rho} \left( \rho \frac{dR}{d\rho} \right) + \frac{\rho^2}{Z} \frac{d^2 Z}{dz^2}= -\frac{1}{F} \frac{d^2 F}{d\phi^2}
        \end{align*}
        逐步分离可得
        \begin{equation*}
            \begin{cases}
                \rho^2R'' + \rho R' - (n^2 + \mu^2\rho^2) R = 0, \\
                F'' + n^2 F = 0, \\
                Z'' + \mu^2 Z = 0
            \end{cases}
        \end{equation*}
        由边界条件可知，$Z(0) = Z(L) = 0$，因此
        \begin{equation*}
            Z_m(z) = \sin\left(\frac{m\pi z}{L}\right)
        \end{equation*}
        本题的边界值都与 \(\phi\) 无关，圆柱区域也具有旋转对称性，因此可以直接取
        \begin{equation*}
            F(\phi) = 1, \qquad n = 0
        \end{equation*}
        因此径向方程可以写为
        \begin{equation*}
            \rho^2R'' + \rho R' - \left(\frac{m\pi}{L}\right)^2\rho^2 R = 0
        \end{equation*}
        这是一个 Bessel 方程，其解为 Bessel 函数，由于函数值在 $\rho =0 $ 处应保持有限，可得
        \begin{equation*}
            R(\rho) = A_m J_0\left(\frac{m\pi}{L} \rho\right)
        \end{equation*}
        解的形式可以写为
        \begin{equation*}
            \Phi(\rho, \phi, z) = \sum_{m=1}^{\infty} A_m J_0\left(\frac{m\pi}{L} \rho\right) \sin\left(\frac{m\pi z}{L}\right)
        \end{equation*}
        由 $\Phi(a, \phi, z) = 1$ 可得
        \begin{equation*}
            A_m = \frac{1}{J_0\left(\frac{m\pi}{L} a\right)}\frac{2(1-(-1)^{m})}{m\pi}
        \end{equation*}
        因此最终的解为
        \begin{equation*}
            \begin{aligned}
                \Phi(\rho, \phi, z) &= \sum_{m=1}^{\infty} \frac{2(1-(-1)^{m})}{m\pi}\frac{J_0\left(\frac{m\pi}{L}\rho\right)}{J_0\left(\frac{m\pi}{L} a\right)} \sin\left(\frac{m\pi z}{L}\right)\\
                &= \sum_{q=0}^{\infty} \frac{4}{(2q+1)\pi}\frac{J_0\left(\frac{(2q+1)\pi}{L}\rho\right)}{J_0\left(\frac{(2q+1)\pi}{L} a\right)} \sin\left(\frac{(2q+1)\pi z}{L}\right)
            \end{aligned}
        \end{equation*}
    \item[(c)] Heat equation
        \begin{align*}
            \Phi_t &= k \nabla^2 \Phi, \quad 0 \leq \rho \leq a, \, 0 \leq \phi < 2\pi, \, t > 0, \\
            \Phi(a, \phi, t) &= 0, \quad
            \Phi(\rho, \phi, 0) = \rho^2 \cos(2\phi).
        \end{align*}
        初始条件与边界条件均与 $z$ 无关，因此解也与 $z$ 无关。令
        \begin{equation*}
            \Phi(\rho, \phi, t) = R(\rho) F(\phi) T(t)
        \end{equation*}
        并代入二维极坐标下的热方程，得
        \begin{equation*}
            \frac{T'}{kT} = 
            \frac{R''}{R} + \frac{1}{\rho}\frac{R'}{R} + \frac{1}{\rho^2}\frac{F''}{F}
        \end{equation*}
        取分离常数为 $-\lambda^2$，再分离角向变量，可得
        \begin{equation*}
            \begin{cases}
                T'+k\lambda^2T=0, \\
                F''+n^2F=0,\\
                \rho^2R''+\rho R'+(\lambda^2\rho^2-n^2)R=0.
            \end{cases}
        \end{equation*}
        角向周期性要求 $n$ 为非负整数，而初始条件的角向部分为
        $\cos(2\phi)$，故只需取
        \begin{equation*}
            n=2,\qquad F(\phi)=\cos(2\phi).
        \end{equation*}
        相应的径向方程为二阶 Bessel 方程
        \begin{equation*}
            \rho^2R''+\rho R'+(\lambda^2\rho^2-4)R=0,
        \end{equation*}
        其通解为
        \begin{equation*}
            R(\rho)=C J_2(\lambda\rho)+D Y_2(\lambda\rho).
        \end{equation*}
        因为区域包含 $\rho=0$，而 $Y_2$ 在原点发散，所以 $D=0$。
        由边界条件 $R(a)=0$ 可得
        \begin{equation*}
            J_2(\lambda a)=0.
        \end{equation*}
        记 $j_{2,m}$ 为 $J_2$ 的第 $m$ 个正零点，则
        \begin{equation*}
            \lambda_m=\frac{j_{2,m}}{a},\qquad
            T_m(t)=\exp\!\left(-\frac{k j_{2,m}^2}{a^2}t\right).
        \end{equation*}
        因此解可展开为
        \begin{equation*}
            \Phi(\rho,\phi,t)=\cos(2\phi)\sum_{m=1}^{\infty}A_m
            J_2\!\left(\frac{j_{2,m}}{a}\rho\right)
            \exp\!\left(-\frac{k j_{2,m}^2}{a^2}t\right).
        \end{equation*}
        由初始条件和 Fourier--Bessel 正交性可得
        \begin{align*}
            A_m
            &=\frac{\displaystyle\int_0^a \rho^3
            J_2\!\left(\frac{j_{2,m}}{a}\rho\right)\,d\rho}
            {\displaystyle\int_0^a \rho
            J_2^2\!\left(\frac{j_{2,m}}{a}\rho\right)\,d\rho}\\
            &=\frac{2a^2}{j_{2,m}J_3(j_{2,m})}.
        \end{align*}
        最终得到
        \begin{equation*}
            \Phi(\rho,\phi,t)=2a^2\cos(2\phi)
            \sum_{m=1}^{\infty}
            \frac{J_2\!\left(\dfrac{j_{2,m}}{a}\rho\right)}
            {j_{2,m}J_3(j_{2,m})}
            \exp\!\left(-\frac{k j_{2,m}^2}{a^2}t\right).
        \end{equation*}
\end{enumerate}

% ===============================================

\end{document}
